\section{Conclusion}
\label{sec:conclusion}

How long an outstation takes to answer is a signature an observer can profile, and the relay
leaking it cannot be modified. This paper presents a timing-obfuscation framework that moves the
problem into the network: a switch in front of the outstation holds the packets carrying the
device's timing and releases each at a deadline computed when the transaction begins. On the read
path, which is READ and the SELECT phase of select-before-operate, both deadlines are computed
from the same instant, the moment the outstation acknowledges the request; the OPERATE case is
anchored to the request instead. Consequently, the interval an observer sees is the operator's
value rather than the device's.

Evaluation on one physical relay shows the released interval concentrating on the configured
value: its interquartile range falls from 2.78 to 0.006~ms and its standard deviation from 2.61
to 0.63~ms. A classifier trained on that interval before deployment drops from 0.652 balanced
accuracy to three-class chance. The framework protects the exchanges its budget can schedule and
forwards the rest untouched.

Three bounds limit the result. An adversary who retrains recovers to 0.651 by adding the
request-to-acknowledgment interval, which the mechanism does not target; the evidence is one
relay, one campaign and a transaction-class evaluation rather than a comparison between device
models; and the switch's own release instants were never captured, so the residual around the
configured value is a difference between two releases rather than a delay we can place inside the
switch. An instrumented build measures an internal blocker-termination interval of about
1.7~$\mu$s, which characterizes the mechanism without decomposing that residual. Evaluating the
framework on outstations from several vendors, and capturing the release instants on both switch
links, is what would attribute it.
